% =====================================================================
%  LaTeX source for tables, equations and the algorithm block
%  Explainable AI-Based Personalized Learning Path Recommendation System
%  IEEEtran conference format, two columns
%
%  Every number is measured and traceable to results/ in
%  github.com/abhay27singh/explainable-learning-path-recommender
% =====================================================================
%
%  PREAMBLE — add these if not already present
%
%  \usepackage{amsmath,amssymb}
%  \usepackage{booktabs}          % \toprule \midrule \bottomrule
%  \usepackage{multirow}
%  \usepackage{algorithm}
%  \usepackage{algpseudocode}
%
%  PAGE BUDGET. Nine tables will not fit a six-page paper. Include the
%  three marked ESSENTIAL; add others as space allows, in the order given.
%
%    ESSENTIAL   Table I   accuracy
%    ESSENTIAL   Table III cold-concept   <- the contribution
%    ESSENTIAL   Table IV  path quality
%    strong      Table II  ablation (sets up Table III; can be folded into prose)
%    strong      Table VII explanation quality
%    if space    Tables V, VI, VIII, IX
%
% =====================================================================


% ---------------------------------------------------------------------
% TABLE I — predictive accuracy                             [ESSENTIAL]
% Six columns: use table* so it spans both columns.
% ---------------------------------------------------------------------
\begin{table*}[t]
\centering
\caption{Comparative predictive performance. Mean $\pm$ standard deviation over
five-fold cross-validation with folds grouped by student, so that no learner appears
in both training and test partitions. $p$-values are paired $t$-tests against the
proposed model, Holm--Bonferroni corrected across the comparison family; $d$ is
Cohen's $d$ for paired samples.}
\label{tab:accuracy}
\begin{tabular}{lccccc}
\toprule
Model & AUC-ROC & RMSE & ECE & $p$ & $d$ \\
\midrule
Majority (no-skill)          & $0.5000 \pm 0.0000$ & $0.4669 \pm 0.0004$ & $0.0037$ & $<0.001$ & --- \\
SAKT \cite{pandey2019sakt}   & $0.9233 \pm 0.0018$ & $0.2992 \pm 0.0017$ & $0.0081$ & $<0.001$ & $26.16$ \\
DKT \cite{piech2015dkt}      & $0.9435 \pm 0.0017$ & $0.2751 \pm 0.0022$ & $0.0102$ & $0.0009$ & $5.67$ \\
GNN-based                    & $0.9502 \pm 0.0014$ & $0.2665 \pm 0.0020$ & $0.0064$ & $0.0220$ & $2.25$ \\
\textbf{Proposed (KG-DKT)}   & $\mathbf{0.9538 \pm 0.0019}$ & $\mathbf{0.2621 \pm 0.0026}$ & $\mathbf{0.0054}$ & --- & --- \\
\midrule
Improvement vs.\ GNN-based   & $+0.0036$ & $-0.0044$ & & & \\
Improvement vs.\ DKT         & $+0.0103$ & $-0.0130$ & & & \\
\bottomrule
\end{tabular}
\end{table*}


% ---------------------------------------------------------------------
% TABLE II — component ablation
% ---------------------------------------------------------------------
\begin{table}[t]
\centering
\caption{Component ablation under the conventional protocol. Neither component
produces a statistically detectable change in predictive accuracy. Section~V-B
explains why this measurement cannot observe the knowledge graph.}
\label{tab:ablation}
\begin{tabular}{lccc}
\toprule
Variant & AUC-ROC & RMSE & $p$ \\
\midrule
Proposed (full)              & $0.9538 \pm 0.0019$ & $0.2621 \pm 0.0026$ & --- \\
\quad w/o knowledge graph    & $0.9539 \pm 0.0021$ & $0.2618 \pm 0.0025$ & $0.77$ \\
\quad w/o temporal attention & $0.9532 \pm 0.0014$ & $0.2624 \pm 0.0024$ & $0.45$ \\
\bottomrule
\end{tabular}
\end{table}


% ---------------------------------------------------------------------
% TABLE III — cold-concept ablation                         [ESSENTIAL]
% This is the paper's contribution. Give it prominence.
% ---------------------------------------------------------------------
\begin{table}[t]
\centering
\caption{Cold-concept ablation. All assessments are withheld for 27 of the 89
assessed concepts and their observed responses are removed from the model input, so
that they behave as genuinely unassessed concepts; evaluation is restricted to those
concepts alone. The graph-conditioned model wins on every fold.}
\label{tab:coldconcept}
\begin{tabular}{lcc}
\toprule
Variant & AUC-ROC & RMSE \\
\midrule
Proposed, with knowledge graph & $\mathbf{0.9004 \pm 0.0059}$ & $\mathbf{0.3359 \pm 0.0059}$ \\
Proposed, without              & $0.8679 \pm 0.0042$ & $0.3611 \pm 0.0113$ \\
\midrule
Difference & \multicolumn{2}{c}{$+0.0325$ \quad $p = 0.0003$ \quad $d = 5.06$} \\
\bottomrule
\end{tabular}
\end{table}


% ---------------------------------------------------------------------
% TABLE IV — learning path quality                          [ESSENTIAL]
% ---------------------------------------------------------------------
\begin{table*}[t]
\centering
\caption{Learning path quality over 906 recommendation decisions from 400 held-out
learners who passed or achieved distinction. Ground truth is the set of concepts each
learner actually engaged with next, rather than a simulation. All planners operate
under the same prerequisite action mask.}
\label{tab:pathquality}
\begin{tabular}{lccccc}
\toprule
Planner & NDCG@5 & Hit@5 & Prec@5 & Prereq.\ violations & Coverage \\
\midrule
Curriculum (syllabus order) & $0.0612 \pm 0.154$ & $0.1832$ & $0.0453$ & $0.0000$ & $0.452$ \\
Popularity                  & $0.1069 \pm 0.199$ & $0.2936$ & $0.0808$ & $0.0000$ & $0.510$ \\
Random (legal moves)        & $0.1300 \pm 0.218$ & $0.3455$ & $0.0863$ & $0.0000$ & $0.722$ \\
Weakest-first               & $0.1344 \pm 0.207$ & $0.3642$ & $0.0956$ & $0.0000$ & $0.515$ \\
\textbf{Proposed (greedy)}  & $\mathbf{0.1532 \pm 0.238}$ & $\mathbf{0.3687}$ & $\mathbf{0.1051}$ & $\mathbf{0.0000}$ & $\mathbf{0.726}$ \\
\bottomrule
\end{tabular}

\vspace{4pt}
\begin{tabular}{lccc}
\multicolumn{4}{c}{\footnotesize Paired comparison of NDCG@5 against the proposed planner (Holm-corrected)}\\
\toprule
Comparison & Difference & $p$ & Cohen's $d$ \\
\midrule
vs.\ curriculum    & $+0.0921$ & $<0.001$  & $0.36$ (small) \\
vs.\ popularity    & $+0.0464$ & $<0.001$  & $0.17$ (negligible) \\
vs.\ random        & $+0.0232$ & $0.0125$  & $0.09$ (negligible) \\
vs.\ weakest-first & $+0.0189$ & $0.0167$  & $0.08$ (negligible) \\
\bottomrule
\end{tabular}
\end{table*}


% ---------------------------------------------------------------------
% TABLE V — calibration
% ---------------------------------------------------------------------
\begin{table}[t]
\centering
\caption{Calibration of predicted mastery. The temperature is fitted on one half of
each fold's held-out predictions and evaluated on the other, so the reported error is
not self-flattering. Mastery values are displayed to learners inside explanations,
which makes their calibration a correctness requirement rather than a refinement.}
\label{tab:calibration}
\begin{tabular}{lc}
\toprule
 & Expected calibration error \\
\midrule
Before temperature scaling & $0.0084$ \\
After temperature scaling  & $\mathbf{0.0054}$ \\
\midrule
Fitted temperature per fold & \multicolumn{1}{c}{$1.022,\ 1.066,\ 1.068,\ 1.024,\ 1.046$} \\
\bottomrule
\end{tabular}
\end{table}


% ---------------------------------------------------------------------
% TABLE VI — cold-start performance
% ---------------------------------------------------------------------
\begin{table}[t]
\centering
\caption{Predictive accuracy as a function of available interaction history. With no
history the model falls back to the learner-profile prior obtained by $k$-means
clustering over first-window behaviour.}
\label{tab:coldstart}
\begin{tabular}{lc}
\toprule
Prior interactions & AUC-ROC \\
\midrule
$0$        & $0.7421 \pm 0.0353$ \\
$1$--$4$   & $0.9449 \pm 0.0086$ \\
$5$--$9$   & $0.9559 \pm 0.0062$ \\
$10$--$19$ & $0.9539 \pm 0.0056$ \\
$\geq 20$  & $0.9514 \pm 0.0023$ \\
\bottomrule
\end{tabular}
\end{table}


% ---------------------------------------------------------------------
% TABLE VII — explanation quality
% ---------------------------------------------------------------------
\begin{table}[t]
\centering
\caption{Objective explanation quality. A rating scale records whether a reader
found an explanation agreeable; these measures record whether it is faithful to the
model. The low sufficiency indicates that generated explanations were accurate but
incomplete, attributing to the named concept only a small share of the benefit they
described.}
\label{tab:explanation}
\begin{tabular}{lcp{4.1cm}}
\toprule
Measure & Value & Definition \\
\midrule
Fidelity    & $0.5318$ & Rank correlation between cited prerequisite deficits and measured counterfactual improvement \\
Sufficiency & $0.0020$ & Share of predicted gain attributable to the concept named \\
Stability   & $0.5310$ & Jaccard overlap of cited gaps under $\pm 0.02$ mastery perturbation \\
\bottomrule
\end{tabular}
\end{table}


% ---------------------------------------------------------------------
% TABLE VIII — dataset and derived graph
% ---------------------------------------------------------------------
\begin{table}[t]
\centering
\caption{Dataset and derived concept graph. OULAD contains no concept annotations,
prerequisite relations or material titles; concepts are defined as (module,
week-of-study) pairs and edges are derived from course structure and mined from
learner behaviour.}
\label{tab:dataset}
\begin{tabular}{lr}
\toprule
Quantity & Value \\
\midrule
Distinct students                & $28{,}785$ \\
Enrolments                       & $32{,}593$ \\
Clickstream events               & $10{,}655{,}280$ \\
Assessment records               & $173{,}912$ \\
Learning materials               & $6{,}364$ \\
\midrule
Concepts derived                 & $237$ \\
\quad prerequisite edges         & $724$ \\
\quad corequisite edges          & $302$ \\
\quad cycles requiring removal   & $0$ \\
\quad concepts with assessments  & $89$ of $237$ \\
\midrule
Interaction sequences            & $27{,}904$ \\
\quad median events per sequence & $169$ \\
Supervised tokens                & $242{,}670$ \\
Label balance (success/failure)  & $67.9\% / 32.1\%$ \\
Placement proxy validation       & $\rho = 0.975$ \\
\bottomrule
\end{tabular}
\end{table}


% ---------------------------------------------------------------------
% TABLE IX — computational cost
% ---------------------------------------------------------------------
\begin{table}[t]
\centering
\caption{Computational cost, measured on the hardware used. Inference runs on CPU;
training used a single free-tier NVIDIA T4.}
\label{tab:efficiency}
\begin{tabular}{lr}
\toprule
Quantity & Value \\
\midrule
Model parameters                     & $466{,}569$ \\
\quad transformer encoder            & $85\%$ \\
\quad graph encoder                  & $6\%$ \\
End-to-end recommendation latency    & ${\approx}160$\,ms \\
Cached mastery lookup                & $7$\,ms \\
Training per configuration           & ${\approx}25$\,min \\
Full ETL from raw data               & ${\approx}2.5$\,min \\
\bottomrule
\end{tabular}
\end{table}


% =====================================================================
%  EQUATIONS for Section III
%  The submitted paper contains none. Reviewer 3: "the presentation of
%  the proposed algorithm is too conceptual, and details of how to
%  operate the proposed algorithm in practice should be clearly
%  elaborated."
% =====================================================================

% --- A. Knowledge-graph concept encoding -----------------------------
Concept representations are produced by a three-layer graph convolution over the
symmetrically normalised adjacency of the prerequisite graph,
%
\begin{equation}
\hat{A} = D^{-1/2}\,(A + I)\,D^{-1/2},
\label{eq:adjacency}
\end{equation}
%
where $A \in \mathbb{R}^{C \times C}$ is the weighted prerequisite adjacency over
$C = 237$ concepts and $D$ its degree matrix. With $H^{(0)} = E$ a learned embedding
table,
%
\begin{equation}
H^{(l+1)} = \mathrm{Dropout}\!\left(\mathrm{ReLU}\!\left(\mathrm{LayerNorm}\!\left(
\hat{A}\,H^{(l)}\,W^{(l)}\right)\right)\right),
\label{eq:gcn}
\end{equation}
%
for $l = 0,1,2$, giving concept embeddings $Z = H^{(3)} \in \mathbb{R}^{C \times 64}$.
$Z$ is recomputed at every forward pass, so gradients propagate into the graph
structure rather than treating it as a fixed input.

% --- B. Interaction embedding ----------------------------------------
For interaction $t$ with concept $c_t$, response
$r_t \in \{\textit{correct}, \textit{incorrect}, \textit{unlabelled}\}$ and kind
$k_t \in \{\textit{context}, \textit{assessment}\}$,
%
\begin{equation}
x_t = W_c Z[c_t] + W_r e(r_t) + W_k e(k_t) + W_f f_t + W_s s,
\label{eq:embedding}
\end{equation}
%
where $s$ are static learner features and
$f_t = [\log(1 + \text{clicks}_t),\ \Delta t_{\text{prev}},\
\Delta t_{\text{concept}},\ \text{day}_t/\text{span}]$ carries the temporal context.

% --- C. Time-aware attention (the forgetting mechanism) ---------------
Forgetting is modelled as an additive attention bias decaying with elapsed time,
%
\begin{equation}
b^{(h)}_{ij} = -\,\mathrm{softplus}(\gamma_h)\,\log\!\left(1 + \Delta_{ij}\right),
\label{eq:decay}
\end{equation}
%
where $\Delta_{ij}$ is the number of days between interactions $i$ and $j$ and
$\gamma_h$ is learned independently for each attention head, allowing one head to
attend to recent activity while another retains a longer memory. Attention is then
%
\begin{equation}
A^{(h)} = \mathrm{softmax}\!\left(
\frac{Q^{(h)} {K^{(h)}}^{\!\top}}{\sqrt{d_k}} + b^{(h)} + M\right),
\label{eq:attention}
\end{equation}
%
with $M$ the combined causal and padding mask. Setting $\gamma_h = 0$ recovers
ordinary scaled dot-product attention and constitutes the temporal ablation reported
in Table~\ref{tab:ablation}.

% --- D. Prediction and masked objective -------------------------------
Writing $h_t$ for the encoder state after interaction $t$, the probability of success
on the next interaction is
%
\begin{equation}
\hat{y}_t = \sigma\!\left(\mathrm{MLP}\!\left(\left[\,h_{t-1}\,;\,Z[c_t]\,\right]\right)\right),
\label{eq:prediction}
\end{equation}
%
conditioning on the identity of concept $c_t$ but never on its response, so no label
information leaks. The objective is binary cross-entropy restricted to assessment
positions $\mathcal{S}$,
%
\begin{equation}
\mathcal{L} = -\frac{1}{|\mathcal{S}|}\sum_{t \in \mathcal{S}}
\Big[\,y_t \log \hat{y}_t + (1-y_t)\log(1-\hat{y}_t)\,\Big].
\label{eq:loss}
\end{equation}
%
Clickstream interactions update the hidden state but contribute no loss term. This is
what makes a sequence model of length 200 trainable on a dataset in which learners
average only 8.7 assessments each.

% --- E. Calibration ---------------------------------------------------
Predicted probabilities are calibrated by temperature scaling, with $T^\star$ fitted
on held-out predictions,
%
\begin{equation}
T^\star = \operatorname*{arg\,min}_{T} \;
\mathrm{BCE}\!\left(\sigma(z/T),\, y\right),
\qquad
\hat{y} = \sigma\!\left(z / T^\star\right).
\label{eq:temperature}
\end{equation}

% --- F. Mastery state -------------------------------------------------
The mastery vector consumed by the planner evaluates the prediction head against every
concept,
%
\begin{equation}
m_t[c] = \sigma\!\left(\mathrm{MLP}\!\left(\left[\,h_t\,;\,Z[c]\,\right]\right)\right),
\qquad c = 1,\dots,C.
\label{eq:mastery}
\end{equation}

% --- G. Decision process ----------------------------------------------
Recommendation is formulated as a Markov decision process. The state concatenates the
mastery vector, a coverage mask of concepts already recommended, the static learner
features and the remaining budget,
%
\begin{equation}
s_t = \left[\,m_t \,\|\, o_t \,\|\, s \,\|\, b_t\,\right] \in \mathbb{R}^{2C + 57}.
\label{eq:state}
\end{equation}
%
An action $a \in \{1,\dots,C\}$ is admissible only if
%
\begin{equation}
m_t[a] < \tau_{\text{mastered}}
\;\wedge\;
\forall p \in \mathrm{pre}(a):\, m_t[p] \geq \tau_{\text{ready}}
\;\wedge\;
a \in \mathcal{E},
\label{eq:mask}
\end{equation}
%
with $\mathcal{E}$ the concepts of the learner's enrolled modules. Encoding the
pedagogical constraint in the admissible set, rather than expecting a learned policy
to discover it, makes an unprepared recommendation impossible by construction; this is
verified empirically at zero violations in Table~\ref{tab:pathquality}. The reward
%
\begin{equation}
r_t = \alpha \sum_{c} \max\!\left(0, \Delta m[c]\right)
+ \beta\,\mathbb{1}[\text{prereq.\ met}]
- \delta\,\mathbb{1}[\text{redundant}]
- \eta
\label{eq:reward}
\end{equation}
%
uses $\alpha = 1.0$, $\beta = 0.2$, $\delta = 0.5$, $\eta = 0.05$.

% --- H. Planner scoring ------------------------------------------------
Candidates are ranked by the predicted improvement in log-odds relative to an equal
period of no study,
%
\begin{equation}
\mathrm{score}(a) = \sum_{c \in \mathcal{E}}
\left[\,
\mathrm{logit}\, m^{(a)}_{t+1}[c] \;-\; \mathrm{logit}\, m^{(\varnothing)}_{t+1}[c]
\,\right].
\label{eq:score}
\end{equation}
%
Two choices here are load-bearing. The baseline $m^{(\varnothing)}_{t+1}$ is an equal
elapsed period without study rather than the present state: studying advances the
clock, and the decay of \eqref{eq:decay} then attenuates every earlier interaction, so
measuring against the present moment makes every available action appear harmful.
Scoring in log-odds rather than probability is necessary because engaged learners
saturate near $0.98$ on most concepts, where probability differences vanish and the
ranking degenerates.

% --- I. Priority backtracking ------------------------------------------
For a recommended concept $c^\star$, each prerequisite ancestor $a$ is scored by its
shortfall, discounted by distance in the graph,
%
\begin{equation}
\pi(a) = \max\!\left(0,\ \tau_{\text{ready}} - m[a]\right)\cdot \rho^{\,d(a, c^\star) - 1},
\label{eq:priority}
\end{equation}
%
with decay $\rho = 0.6$ and traversal depth at most four. Ancestors are reported in
descending $\pi$, so that an unmet prerequisite immediately preceding the
recommendation outranks an equal shortfall further back.


% =====================================================================
%  ALGORITHM BLOCK
%  Reviewer 3 asked for operational detail. This supplies it.
% =====================================================================
\begin{algorithm}[t]
\caption{Explained recommendation}
\label{alg:recommend}
\begin{algorithmic}[1]
\Require learner history $h$, graph $G$, trained parameters $\theta$, budget $k$
\Ensure $k$ recommendations, each with an explanation
\State $Z \gets \textproc{GCN}(\hat{A})$ \Comment{Eq.~\eqref{eq:gcn}}
\State $m \gets \textproc{Mastery}(h, Z, \theta)$ \Comment{Eq.~\eqref{eq:mastery}, calibrated}
\State $\tau \gets \textproc{Percentiles}(m|_{\mathcal{E}})$ \Comment{learner-relative}
\State $L \gets \{\, a : \textproc{Admissible}(a, m, \tau, G) \,\}$ \Comment{Eq.~\eqref{eq:mask}}
\If{$L = \varnothing$}
  \State $L \gets \textproc{Relax}(m, \tau)$ \Comment{report relaxation in the explanation}
\EndIf
\State $m^{(\varnothing)} \gets \textproc{Mastery}(\textproc{Idle}(h), Z, \theta)$
\ForAll{$a \in L$} \Comment{one batched forward pass}
  \State $m^{(a)} \gets \textproc{Mastery}(h \oplus (a, \text{correct}), Z, \theta)$
  \State $\mathrm{score}(a) \gets \sum_{c \in \mathcal{E}} \left[\mathrm{logit}\,m^{(a)}[c] - \mathrm{logit}\,m^{(\varnothing)}[c]\right]$
\EndFor
\State $A \gets \textproc{TopK}(L, \mathrm{score}, k)$
\ForAll{$a \in A$}
  \State $\mathcal{G} \gets \textproc{Backtrack}(G, m, a, \tau)$ \Comment{Eq.~\eqref{eq:priority}}
  \State $\mathcal{C} \gets \textproc{Counterfactual}(a, \mathcal{G}, \theta)$ \Comment{computed, not asserted}
  \State $\textproc{Render}(a, \mathcal{G}, \mathcal{C})$
\EndFor
\State \Return $A$ with explanations
\end{algorithmic}
\end{algorithm}


% =====================================================================
%  BIBLIOGRAPHY ENTRIES referenced above
% =====================================================================
%
% \bibitem{piech2015dkt}
% C.~Piech, J.~Bassen, J.~Huang, S.~Ganguli, M.~Sahami, L.~Guibas, and
% J.~Sohl-Dickstein, ``Deep knowledge tracing,'' in \emph{Advances in Neural
% Information Processing Systems}, vol.~28, 2015, pp.~505--513.
%
% \bibitem{pandey2019sakt}
% S.~Pandey and G.~Karypis, ``A self-attentive model for knowledge tracing,'' in
% \emph{Proc. 12th Int. Conf. Educational Data Mining (EDM)}, 2019, pp.~384--389.
%
% \bibitem{kuzilek2017oulad}
% J.~Kuzilek, M.~Hlosta, and Z.~Zdrahal, ``Open University Learning Analytics
% dataset,'' \emph{Scientific Data}, vol.~4, Art. no. 170171, 2017,
% doi: 10.1038/sdata.2017.171.
%
% \bibitem{guo2017calibration}
% C.~Guo, G.~Pleiss, Y.~Sun, and K.~Q. Weinberger, ``On calibration of modern neural
% networks,'' in \emph{Proc. 34th Int. Conf. Machine Learning (ICML)}, 2017,
% pp.~1321--1330.
%
% \bibitem{holm1979}
% S.~Holm, ``A simple sequentially rejective multiple test procedure,''
% \emph{Scandinavian Journal of Statistics}, vol.~6, no.~2, pp.~65--70, 1979.
%
% NOTE. Two citations are required by the methods above and are absent from the
% current reference list: Guo et al. for temperature scaling, and Holm for the
% multiple-comparison correction. Pandey & Karypis is required for the SAKT baseline.
